\documentclass{article}
\usepackage{mathtools} 
\usepackage{fontspec}
\usepackage[UTF8]{ctex}
\usepackage{amsthm}
\usepackage{mdframed}
\usepackage{xcolor}
\usepackage{amssymb}
\usepackage{amsmath}


% 定义新的带灰色背景的说明环境 zremark
\newmdtheoremenv[
  backgroundcolor=gray!10,
  % 边框与背景一致，边框线会消失
  linecolor=gray!10
]{zremark}{说明}

% 通用矩阵命令: \flexmatrix{矩阵名}{元素符号}{行数}{列数}
\newcommand{\flexmatrix}[4]{
  \[
  #1 = \begin{pmatrix}
    #2_{11}     & #2_{12}     & \cdots & #2_{1#4}   \\
    #2_{21}     & #2_{22}     & \cdots & #2_{2#4}   \\
    \vdots      & \vdots      & \ddots & \vdots     \\
    #2_{#31}    & #2_{#32}    & \cdots & #2_{#3#4}
  \end{pmatrix}
  \]
}

% 简化版命令（默认矩阵名为A，元素符号为a）: \quickmatrix{行数}{列数}
\newcommand{\quickmatrix}[2]{\flexmatrix{A}{a}{#1}{#2}}

\begin{document}
\title{注释 15.2}
\author{张志聪}
\maketitle

\begin{zremark}
  傅里叶级数从$2\pi$推广到$2l$过程中使用了以下命题：

  变量代换的可逆性
\end{zremark}

举例说明这个问题：

设函数 $f: \mathbb{R} \to \mathbb{R}$，常数 $a \neq 0$。定义映射 $\phi: \mathbb{R} \to \mathbb{R}$ 为 $\phi(t) = a t$。

\noindent\textbf{第一步：} 令 $x = \phi(t) = a t$，则复合函数
\[
  F(t) = (f \circ \phi)(t) = f(\phi(t)) = f(a t).
\]

\noindent\textbf{第二步：} 映射 $\phi$ 是可逆的，其逆映射为 $\phi^{-1}(x) = \dfrac{x}{a}$，
令 $t = \phi^{-1}(x) = \dfrac{x}{a}$。代入 $F(t)$ 得
\[
  F(\phi^{-1}(x)) = (f \circ \phi)(\phi^{-1}(x)) = f(\phi(\phi^{-1}(x))) = f(x).
\]

即
\[
  f\left( a \cdot \frac{x}{a} \right) = f(x).
\]

\end{document}